\chapter{CONCLUSIONS AND RECOMMENDATIONS}
\section{Conclusions}
From my construction-site experience in Pokhara, I considered waterproofing details, brickwork junctions and plaster defects to be recurring workmanship concerns. The Lamachur residential-house photographs illustrate work associated with a door-threshold waterproofing correction. Other photographs show roof, masonry and road-concrete conditions at separate locations. These examples support a qualitative engineering discussion but do not establish how frequently rework occurred across Pokhara or its total financial impact.

The technical analysis shows why each problem requires more than a surface-level explanation. For waterproofing, a complete assembly depends on appropriate substrate, membrane continuity, drainage and compatible transitions rather than tile or grout alone. For PCC pavement, cracking must be related to concrete behaviour, loads, support, curing and joint design; the absence of distributed steel does not automatically imply noncompliance in jointed plain concrete. For masonry and plaster, weak interfaces, inappropriate moisture preparation or missing applicable details may contribute to cracking, but no single field observation is sufficient to establish its cause.

Across the three categories, an important recurring quality-management requirement is verification at concealed or difficult-to-rework stages. Inspection before placing finishes, before pavement acceptance, and before plastering can help identify discrepancies while correction is less disruptive. The general construction-rework literature likewise supports the importance of good records and early control \parencite{forcada2017,safapour2019,love2022}. Nevertheless, the supplied material does not establish a frequency ranking, measured monetary loss or completion delay for the selected projects. Such results must await identifiable project records and verified corrective interventions.

\section{Recommendations for Practice}
\subsection{Waterproofing}
Use an approved, compatible wet-area or roof waterproofing assembly and identify substrate, drainage and termination requirements before installation. Check continuity at floor--wall junctions, drains, pipes, door thresholds, service penetrations and repairs to temporary support holes. Confirm acceptable falls and outlet conditions. Apply manufacturer's specified preparation, application, curing and testing instructions; record acceptance before tiling or covering. Use spacers and properly dimensioned joints where specified, but never treat grout alone as the designed waterproof barrier.

\subsection{Concrete pavement}
For PCC pavement associated with a building project, compare the constructed slab with approved load, thickness, concrete, subbase, joint and curing requirements. Establish whether the intended system is plain, jointed, reinforced or continuously reinforced before criticizing the absence of distributed bars. Check joint placement and timing, support preparation, finishing and curing records. Where cracking occurs, map its extent and assess function before choosing monitoring, sealing, local repair or replacement. Only verified correction of completed work should be reported as rework.

\subsection{Plaster and masonry}
Inspect the junction detail between intersecting masonry walls against approved drawings; use bonding or engineered connectors appropriate to the specified wall system. Verify whether lintel or sill bands are required by the actual approved design and applicable Nepal building-code provisions. Assess brick suction and the correct moisture-conditioning procedure, and control mortar curing and timing of plaster application. Map cracks and investigate their origin before undertaking patch repairs that may conceal an unresolved movement mechanism.

\section{Recommendations for Research Documentation}
A researcher wishing to move beyond this qualitative stage should confirm the contributor's identity and the original-language narrative, then relate each potential event to an anonymised project, time period, construction element, original work, observed condition and actual corrective action. Approved drawings, inspection results, dated photographs, site diaries, measured quantities and schedule records should be linked to each case. A separate assessment should establish the cause and the directly attributable impacts rather than presuming either from a photograph or reported workmanship concern.

The PCC portion requires particular attention to project eligibility: an internal access pavement forming part of an RCC-building work package can be considered, whereas an independent municipal road cannot be included without an approved change in study scope. Cross-case frequency, cost and delay comparisons should be undertaken only after the relevant project denominator, admission decisions and record completeness are documented.

\section{Limitations and Further Work}
This report is based on the site observations I described, the selected construction photographs and technical literature. For the photographed work, not all original drawings, daily records and measurements have yet been cross-checked in this report; consequently, I have not presented the unpublished repair-cost figures, causal rankings or claimed overall completion delays as verified results. Therefore it offers an engineering interpretation of reported conditions and a concrete verification procedure, not a completed numerical impact study. Further work should reconcile existing original diaries, drawings and repair records, then test alternative mechanisms and document actual corrective works before drawing stronger conclusions about construction-rework causes and impacts within Pokhara.
